\documentclass[11pt,a4paper]{article}
\usepackage[T1]{fontenc}
\usepackage[utf8]{inputenc}
\usepackage{libertinus}
\usepackage[scaled=.94]{sourcesanspro}
\usepackage{microtype}
\usepackage[a4paper,top=23mm,bottom=24mm,left=25mm,right=25mm,headheight=15pt]{geometry}
\usepackage{parskip}
\usepackage{multicol}
\setlength{\parindent}{0pt}
\usepackage{xcolor}
\definecolor{hrgreen}{HTML}{7DB87A}
\definecolor{hrrust}{HTML}{B0603A}
\definecolor{hrink}{HTML}{202622}
\definecolor{hrmuted}{HTML}{667068}
\definecolor{hrsoft}{HTML}{F4F6F2}
\usepackage{graphicx}
\usepackage{booktabs,longtable,tabularx,array}
\usepackage{enumitem}
\setlist{itemsep=2pt,topsep=5pt}
\usepackage{amsmath,amssymb,mathtools}
\usepackage{tikz}
\usetikzlibrary{positioning,arrows.meta,calc,fit}
\usepackage[most]{tcolorbox}
\tcbset{camelotbox/.style={enhanced,breakable,colback=hrsoft,colframe=hrgreen!70!black,boxrule=.7pt,arc=2pt,left=3mm,right=3mm,top=2.5mm,bottom=2.5mm,fonttitle=\sffamily\bfseries,coltitle=hrink}}
\usepackage{titlesec}
\titleformat{\section}{\Large\bfseries\sffamily\color{hrink}}{\thesection}{.7em}{}
\titleformat{\subsection}{\large\bfseries\sffamily\color{hrink}}{\thesubsection}{.7em}{}
\titleformat{\subsubsection}{\normalsize\bfseries\sffamily\color{hrink}}{\thesubsubsection}{.7em}{}
\titlespacing*{\section}{0pt}{1.8em}{.65em}
\titlespacing*{\subsection}{0pt}{1.3em}{.45em}
\usepackage{fancyhdr}
\pagestyle{fancy}
\fancyhf{}
\fancyhead[L]{\sffamily\footnotesize Where Does the Function Go?}
\fancyhead[R]{\sffamily\footnotesize Hedegreen Research · MI-001}
\fancyfoot[C]{\sffamily\footnotesize \thepage}
\renewcommand{\headrulewidth}{0.25pt}
\usepackage{xurl}
\usepackage[colorlinks=true,linkcolor=hrgreen!45!black,citecolor=hrrust!70!black,urlcolor=hrrust!70!black]{hyperref}
\hypersetup{pdftitle={Where Does the Function Go?},pdfauthor={Dennis Hedegreen},pdfsubject={Mechanistic interpretability research proposal},pdfkeywords={mechanistic interpretability, H-Neurons, J-space, functional reconstitution, recoverability}}
\usepackage{bookmark}
\usepackage{listings}
\lstset{basicstyle=\ttfamily\footnotesize,breaklines=true,columns=fullflexible,frame=single,rulecolor=\color{hrgreen!40},backgroundcolor=\color{hrsoft},showstringspaces=false}
\usepackage{fvextra}
\usepackage{csquotes}
\usepackage[backend=biber,style=authoryear,maxcitenames=2,maxbibnames=99,doi=true,url=true,isbn=false]{biblatex}
\addbibresource{references.bib}
\setlength{\bibitemsep}{0.55\baselineskip}
\usepackage{caption}
\captionsetup{font=small,labelfont={bf,sf},textfont=it}
\usepackage{etoolbox}
\AtBeginEnvironment{longtable}{\small}
\renewcommand{\arraystretch}{1.15}
\providecommand{\tightlist}{%
  \setlength{\itemsep}{0pt}\setlength{\parskip}{0pt}}
\providecommand{\passthrough}[1]{#1}
\newcommand{\statuspill}[1]{\tcbox[on line,colback=hrgreen!10,colframe=hrgreen!60!black,boxrule=.5pt,arc=1.5pt,left=1mm,right=1mm,top=.4mm,bottom=.4mm]{\sffamily\scriptsize\bfseries #1}}
\newcommand{\metarow}[2]{\textsf{\footnotesize\textbf{#1}} & \footnotesize #2 \\\addlinespace[1.5pt]}
\begin{document}

% Cover
\begin{titlepage}
\color{hrink}
\vspace*{5mm}
{\sffamily\small\bfseries\color{hrgreen!55!black}\MakeUppercase{Hedegreen Research / Camelot}}\hfill{\sffamily\small\color{hrmuted} MODEL INTERNALS · MI-001}
\vspace{24mm}

{\sffamily\fontsize{31}{34}\selectfont\bfseries\color{hrink} Where Does the Function Go?\par}
\vspace{5mm}
{\fontsize{15}{19}\selectfont\color{hrink} A Research Proposal for Functional Reconstitution and\\Pre-Commitment Recoverability in Language Models\par}
\vspace{8mm}
{\color{hrgreen}\rule{\textwidth}{1.2pt}}
\vspace{11mm}

{\large\itshape\color{hrink} A detectable signal is not automatically the mechanism that produced the behaviour.\par}
\vfill
\centering
\input{figures/cover_flow.tex}
\vfill
\raggedright
\begin{tabular}{@{}p{31mm}p{110mm}@{}}
\metarow{Author}{Dennis Hedegreen}
\metarow{Publisher}{Hedegreen Research}
\metarow{Date}{21 September 2026}
\metarow{Issue}{MI-001}
\metarow{Version}{v0.8 release candidate}
\metarow{DOI}{Pending}
\metarow{Contact}{contact@hedegreenresearch.com}
\metarow{Status}{Working paper — research proposal / pre-results}
\metarow{Companion}{\href{https://hedegreenresearch.com/articles/where-does-the-function-go/}{Article 200: \emph{Where Does the Function Go?}}}
\end{tabular}
\end{titlepage}
\color{hrink}
\pagenumbering{arabic}
\setcounter{page}{1}

\section*{Abstract}
Mechanistic interpretability frequently identifies internal model components whose activity predicts or causally affects a target behaviour. Such localization results are informative, but they do not by themselves establish that the identified component is the unique or necessary implementation of the function. Large language models may contain redundant routes, distributed representations, or mechanisms that reorganize under further optimization.

This proposal develops two connected experimental programmes.

The first tests \textbf{functional reconstitution}. Starting from hallucination-associated H-Neurons reported by Gao et al.~(2025), the experiment asks whether a behavioural function that is selectively disrupted by persistent neuron-level constraints can recover during continued optimization through a new internal implementation. Version 0.6 strengthens the redundancy test further: candidate backups are measured conditionally at time zero, before adaptation, and later outcomes are separated into pre-existing self-repair, adaptive amplification of a weak backup, distributed redistribution, restricted-subspace compensation, and de novo or near-null reconstitution. The strongest label is reserved for recovery that cannot already be explained by a conditionally causal backup at time zero.

The second programme tests \textbf{pre-commitment recoverability}. Rather than asking whether a model ``knows'' that an answer is wrong, the proposal asks whether hidden state before commitment predicts whether the model's own continuation distribution still assigns meaningful probability to a verified correction. Natural-support continuation with termination allowed remains primary, but version 0.6 adds a one-token termination-barrier control so stopping policy can be separated from post-stop recoverability. Forced-compute continuation with termination constrained is retained only as a separate secondary intervention. Recoverability is compared against simpler confidence, prompt-only and residual-stream baselines. J-space representations exposed by the Jacobian lens remain one candidate explanatory level, not a privileged assumption.

A later bridge experiment tests whether local H-Neuron-associated mechanisms and distributed J-space-like states are causally coupled. A smaller proposition-substitution task is retained as a controlled model organism for self-monitoring failures. Version 0.7 also adds an optional downstream state-transfer branch inspired by Cache-to-Cache communication: if a source model fails, can a separate receiver extract task-relevant information from the source model's KV-cache that is not recovered from the same visible text alone? This branch is explicitly separated from the core recoverability claim.

The proposal is intentionally pre-results. It defines targets, controls, decision rules, failure modes and compute gates before substantive adaptation runs are conducted.

\vspace{3mm}
\textbf{Keywords:} mechanistic interpretability; large language models; neuron localization; H-Neurons; J-space; causal intervention; model adaptation; recoverability.

\begin{tcolorbox}[camelotbox,title={Preferred citation}]
Hedegreen, Dennis. 2026. \emph{Where Does the Function Go? A Research Proposal for Functional Reconstitution and Pre-Commitment Recoverability in Language Models}. Hedegreen Research Working Paper, MI-001, v0.8 release candidate. DOI pending.
\end{tcolorbox}

\begin{tcolorbox}[camelotbox,title={Companion article}]
Public-facing companion: \href{https://hedegreenresearch.com/articles/where-does-the-function-go/}{Article 200, \emph{Where Does the Function Go?}}. The article gives the readable research narrative; this working paper carries the current protocol, controls, decision rules and implementation boundaries.
\end{tcolorbox}

\begin{tcolorbox}[camelotbox,title={Protocol revision note -- v0.8}]
Version 0.8 corrects only the layout of Figure 1. The v0.7 research protocol and all scientific claims are unchanged.


This release preserves the v0.6 adversarial-review revisions and adds Cache-to-Cache (C2C) as a bounded downstream state-transfer test. The new branch asks whether a failed model's internal KV-cache can improve a separate receiver beyond the same visible text, while treating the cache fuser as a potential confound that must be trained and frozen on disjoint data. C2C is not used as evidence for J-space, recoverability, consciousness, or self-knowledge. No empirical result is added by this revision.
\end{tcolorbox}

\clearpage
\section*{How to read this paper}
\addcontentsline{toc}{section}{How to read this paper}
This is a pre-results research proposal. It does not claim that functional reconstitution occurs, that H-Neurons form a universal hallucination circuit, that recoverability is equivalent to ``knowing'' an answer is wrong, or that J-space establishes consciousness. The paper specifies experiments, gates, controls and implementation requirements that must survive before stronger claims are available.

\section*{What this proposal tests}
\addcontentsline{toc}{section}{What this proposal tests}
The proposal asks two primary empirical questions. First, if a selectively causal internal locus is kept persistently unavailable while optimization continues, can the affected behaviour return through a new selectively causal implementation? Second, before a model commits to an incorrect or incomplete answer, does its hidden state predict whether bounded additional computation can rescue the trajectory beyond simpler confidence signals? A later bridge experiment tests whether local H-Neuron-associated mechanisms and distributed J-space-like states are causally related.

\clearpage
\section*{The Proposal in Two Minutes}
\addcontentsline{toc}{section}{The Proposal in Two Minutes}
\begin{multicols}{2}
\textbf{The localization problem.} A neuron or direction can predict a behaviour without being the unique mechanism that implements it. The proposal separates predictive localization, causal contribution, functional necessity, and implementation identity.

\textbf{Functional reconstitution.} H-Neurons are first replicated and causally tested. If a selective locus survives Gate 0, it is kept persistently unavailable while the rest of the model is allowed to adapt. A strong reconstitution claim requires R0--R4.

\textbf{Redundancy before learning.} Alternative routes are tested conditionally immediately after masking and before adaptation. A strong de novo-reconstitution claim requires that later recovery cannot already be explained by a conditionally causal backup at time zero. Weak backups that grow during learning are classified separately as adaptive amplification.

\columnbreak

\textbf{Recoverability, not ``knowing''.} The primary question is whether the model's own continuation distribution, from the exact pre-stop prefix with termination still allowed, contains probability mass leading to verified correction. Forced continuation is a separate secondary intervention, not the definition of recoverability.

\textbf{J-space is a candidate.} J-space is tested as one distributed explanatory level. It receives no privilege unless it adds predictive or causal value beyond simpler residual-stream and confidence features.

\textbf{Compute comes last.} Replication, mask integrity, time-zero mapping, recoverability instrumentation, and external methodological attack come before substantive constrained adaptation. The proposal asks for the cheapest experiment capable of falsifying the programme.
\end{multicols}

\clearpage
\tableofcontents
\clearpage

\section{Research Objective}\label{research-objective}

The central methodological question is:

\begin{quote}
\textbf{When a behaviour can be localized inside a language model, what exactly has been localized, and what happens to the function when that location is made persistently unavailable?}
\end{quote}

The proposal separates four claims that should not be silently collapsed:

\begin{enumerate}
\def\labelenumi{\arabic{enumi}.}
\tightlist
\item
  \textbf{Predictive localization} --- an internal variable predicts a behaviour.
\item
  \textbf{Causal contribution} --- intervening on that variable changes the behaviour.
\item
  \textbf{Functional necessity} --- the behaviour depends selectively on that variable under the tested conditions.
\item
  \textbf{Implementation identity} --- the variable is the mechanism, rather than one implementation, readout, bottleneck, or correlated component of a larger computation.
\end{enumerate}

The experiments are designed to move from weaker to stronger claims only when each preceding gate survives.

A second question concerns internal state before output commitment:

\begin{quote}
\textbf{Does pre-commitment hidden state predict whether additional computation can rescue the current trajectory, beyond ordinary confidence signals?}
\end{quote}

This proposal calls that property \textbf{recoverability}.

The final programme-level question is whether the two levels are related:

\begin{quote}
\textbf{Are sparse local mechanisms associated with answer commitment causally coupled to a distributed unresolved-state representation?}
\end{quote}

\begin{center}\rule{0.5\linewidth}{0.5pt}\end{center}

\section{Scope and Non-Claims}\label{scope-and-non-claims}

This working paper proposes experiments. It does not report that any of the following have been demonstrated:

\begin{itemize}
\tightlist
\item
  that H-Neurons are a universal hallucination circuit;
\item
  that hallucination is structurally necessary to next-token prediction;
\item
  that a function will reconstitute after intervention;
\item
  that J-space is the unique or correct representation of recoverability;
\item
  that pre-commitment recoverability is equivalent to ``knowing one is wrong'';
\item
  that any proposed state is evidence of consciousness;
\item
  that neuron-level localization is sufficient to identify a complete mechanism.
\end{itemize}

The earlier Hedegreen Research proposal \emph{Compliance Is All You Learned}~\parencite{hedegreen2026compliance} used iterative inference-time replacement while model weights remained fixed. That design remains useful as provenance and as a source of controls, but it is \textbf{superseded for the reconstitution claim} because a fixed-weight model cannot adaptively rebuild a mechanism.

\begin{center}\rule{0.5\linewidth}{0.5pt}\end{center}

\section{Background and Motivation}\label{background-and-motivation}

\subsection{H-Neurons}\label{h-neurons}

\textcite{gao2025hneuron} report a sparse subset of MLP intermediate neurons whose contribution profiles predict hallucination across tested model families and tasks. Their experiments also connect H-Neuron interventions to several forms of over-compliance, including invalid-premise continuation, misleading-context acceptance, sycophancy, and harmful-instruction compliance. The earlier Hedegreen Research proposal adopted the same Llama-3.1-8B-Instruct target to preserve comparability with that pipeline.

The original proposal also retained several useful methodological components:

\begin{itemize}
\tightlist
\item
  replication before novel intervention;
\item
  CETT-based contribution profiling;
\item
  threshold sensitivity;
\item
  held-out prediction;
\item
  sham-hook controls;
\item
  layer/location-matched non-H controls;
\item
  random ablation;
\item
  fluency/perplexity monitoring;
\item
  fixed prompt sets and random seeds;
\item
  explicit outcome definitions.
\end{itemize}

Those components remain useful. The interpretation changes.

\subsection{Localization is not mechanism identity}\label{localization-is-not-mechanism-identity}

The broader model-editing literature warns against equating localization with mechanism identity. \textcite{hase2023localization} showed that causal localization did not reliably identify the most effective factual-editing locations. Later work strengthens the warning that successful behavioural editing at an internal location is not itself sufficient evidence that the behaviour is encoded there \autocite{wang2025editing}.

A second problem is \textbf{conditional self-repair}. A component can appear unimportant in the intact model because a dormant backup only becomes necessary after a primary route is removed. \textcite{gong2026coax} formalize this as conditional circuit completion: the relevant causal quantity is how much a remaining component's ablation effect grows once the primary set is absent. This is directly relevant to the present redundancy-versus-reconstitution problem.

The present proposal therefore requires both \textbf{selectivity} and \textbf{conditional training-time dynamics}. Ordinary marginal importance in the intact model is insufficient.

\subsection{J-space and distributed workspace-like representations}\label{j-space-and-distributed-workspace-like-representations}

\textcite{gurnee2026workspace} introduce the Jacobian lens (J-lens), which maps earlier residual-stream activations toward final-layer representational coordinates and decodes them using the model's output vocabulary. Their J-space construction uses sparse nonnegative combinations of an overcomplete set of J-lens vectors and is proposed as a functionally privileged representational component with several workspace-like properties.

This proposal treats J-space as a candidate explanatory level. It must compete against simpler baselines.

\subsection{Hidden correctness and continuation signals are prior art}\label{hidden-correctness-signals-are-prior-art}

Recent work already reports that hidden states can contain information predictive of future correctness before answer generation \autocite{cencerrado2025noanswer,luo2026missing,dicicco2026code}. The present proposal therefore does \textbf{not} claim novelty for internal correctness prediction.

Continuation itself can reveal information that forced extraction misses. \textcite{wang2026detection} report a detection--extraction gap in reasoning models: free continuations from partial prefixes can recover correct answers in cases where forced extraction fails. That result strengthens the need to distinguish natural continuation from an intervention that forces additional computation.

The narrower target here is whether pre-commitment state predicts \textbf{rescue under a precisely defined continuation regime}, and whether the internal signal adds information beyond prompt difficulty, confidence and generic residual-stream features.

\subsection{KV-cache as a communication medium}\label{kv-cache-communication}

\textcite{fu2026c2c} introduce Cache-to-Cache (C2C), a direct communication method in which a learned module projects and fuses one model's KV-cache into another model's cache. Their motivating oracle experiment is especially relevant here: they prefill a model on exemplars plus a question, discard the exemplar-cache segment, retain only the question-aligned cache, and still observe higher accuracy than direct question-only decoding at the same cache length. This supports a limited claim that the internal cache representation of a visible token sequence can carry task-relevant semantic structure that is not exhausted by token count alone.

C2C then trains a cache fuser while freezing the communicating language models. Across several model families and sizes, the authors report gains over receiver-only and text-to-text collaboration, while also showing that poor source states or poorly chosen layers can hurt performance. The result therefore does \textbf{not} imply that all internal state is useful, that KV-cache is a global workspace, or that a failed model ``knows'' the answer. It establishes a technically concrete route by which one model's internal state can alter another model's behaviour without first being compressed into intermediate text.

For the present programme, C2C motivates a downstream test of \textbf{representational surplus}: whether a failed source trajectory contains receiver-useful information in its internal cache beyond what a receiver obtains from the same visible text. Because the C2C fuser itself is learned, any such experiment must separate source-state information from information injected by the fuser's training distribution. It is therefore an optional later branch, not evidence for the core recoverability construct.

\begin{center}\rule{0.5\linewidth}{0.5pt}\end{center}

\section{Operational Definitions}\label{operational-definitions}

\subsection{Target phenotype}\label{target-phenotype}

A \textbf{target phenotype} is a behaviour measured by a pre-specified evaluation suite.

For the pilot, the primary causal phenotype is \textbf{FaithEval --- Counterfactual Context compliance rate}.

Rationale:

\begin{itemize}
\tightlist
\item
  it directly tests over-compliance with misleading context;
\item
  Gao et al.~use greedy decoding with a maximum of 256 new tokens;
\item
  the primary outcome can be extracted with a rule-based parser against dataset labels, avoiding an external LLM judge for the main endpoint.
\end{itemize}

Pilot secondary phenotypes:

\begin{itemize}
\tightlist
\item
  FalseQA invalid-premise compliance;
\item
  Sycophancy answer-flip rate.
\end{itemize}

Jailbreak is not a primary pilot phenotype. It remains a later replication/robustness target because its attack-success endpoint introduces an additional safety-evaluator layer and alignment-specific confounds.

The H-Neuron \textbf{identification} stage remains distinct from the behavioural phenotype stage:

\begin{itemize}
\tightlist
\item
  identification source: TriviaQA faithful vs hallucinated responses;
\item
  OOD diagnostic sets: NQ-Open and BioASQ;
\item
  NonExist retained as a diagnostic rather than a pass/fail gate for Llama-3.1-8B because the source paper reports weak performance for that model on this particular set.
\end{itemize}

\subsection{Candidate locus}\label{candidate-locus}

A \textbf{candidate locus} is an internal neuron, neuron population, direction, circuit, or representational component selected by a pre-specified localization procedure.

A candidate locus is not assumed to be the complete mechanism.

\subsection{Selective causal contribution}\label{selective-causal-contribution}

A locus demonstrates \textbf{selective causal contribution} when intervention on that locus changes the target phenotype more strongly than matched control interventions while preserving sufficient general capability to rule out gross model damage as the explanation.

\subsection{Persistent constraint}\label{persistent-constraint}

A \textbf{persistent constraint} is an intervention kept active on every relevant forward pass during adaptation.

For the first H-Neuron experiment, the candidate implementation is a fixed mask on selected MLP intermediate dimensions applied after the gated activation and before the down-projection.

The implementation must be audited to confirm that the constrained dimensions remain unavailable throughout training.

\subsection{Functional reconstitution and time-zero backup state}\label{functional-reconstitution}

R0 establishes \textbf{selective causal contribution}, not mechanism identity. Stronger claims require additional gates and a conditional time-zero backup audit.

The core R0--R4 sequence remains:

\begin{itemize}
\tightlist
\item \textbf{R0 --- Selective initial effect:} the original target has a selective causal effect beyond matched controls.
\item \textbf{R1 --- Phenotype loss:} persistent constraint measurably reduces the target phenotype before adaptation.
\item \textbf{R2 --- Training-dependent recovery:} the phenotype recovers during adaptation relative to masked no-training and other relevant controls.
\item \textbf{R3 --- Training-dependent replacement change:} a candidate replacement becomes more predictive and/or more conditionally causal across training checkpoints than it was at time zero.
\item \textbf{R4 --- Selective causal dependence:} intervention on the candidate replacement selectively damages the recovered phenotype beyond matched control interventions.
\end{itemize}

R0--R4 are interpreted together with a separate \textbf{B0 conditional backup state} measured immediately after the original locus is constrained and before adaptation:

\begin{itemize}
\tightlist
\item \textbf{B0-A --- Near-null backup:} no candidate alternative carries more than a pre-registered practically negligible conditional causal effect.
\item \textbf{B0-B --- Weak backup:} a candidate route has non-negligible but limited conditional causal effect at time zero.
\item \textbf{B0-C --- Strong pre-existing backup:} an alternative route already carries substantial conditional causal effect before adaptation.
\end{itemize}

The strongest label, \textbf{de novo / near-null functional reconstitution}, is reserved for B0-A plus R0--R4 and cross-seed replication. B0-B followed by a large training-dependent increase is classified as \textbf{adaptive amplification of pre-existing redundancy}. B0-C is classified as \textbf{self-repair / pre-existing backup}, not reconstitution.

Version 0.6 freezes the \emph{rule} for the B0-A practical-equivalence margin before adaptation. Let $\Delta_{R1}$ be the absolute pre-adaptation phenotype loss caused by the persistent H-mask. Define

\[
\delta_{B0}=\min(0.02,\;0.20\,|\Delta_{R1}|),
\]

where effects are expressed as absolute proportions. Because Gate 0 already requires at least a five-percentage-point phenotype shift, this rule yields a candidate equivalence margin between one and two percentage points. B0-A is available only when a paired 90\% confidence interval for the candidate backup's conditional causal risk difference lies wholly inside $[-\delta_{B0},+\delta_{B0}]$. This is the paired-binary analogue of an equivalence test; a generic two-sample $t$-test is not used for binary compliance outcomes.

The strongest label is also \textbf{power-gated}: if the fixed evaluation set cannot provide at least 80\% power for the frozen equivalence margin under the pre-adaptation paired design, B0-A remains unresolved and de novo / near-null reconstitution cannot be claimed. The evaluation set must then be enlarged or the strongest label remains unavailable. The margin rule may not be changed after any adaptation result is observed.

\subsection{Compensation}\label{compensation}

\textbf{Compensation} is behavioural recovery through an adaptation mechanism whose allowed parameter space or architecture materially predetermines where replacement may occur, without sufficient evidence for broader network reorganization.

A LoRA-only recovery may therefore qualify as compensation rather than full reconstitution unless the interpretation is explicitly restricted to the permitted adaptation subspace.

\subsection{Recoverability}\label{recoverability}

Version 0.5 separates three intervention-relative constructs rather than treating hard stop suppression as the definition of recoverability.

\textbf{Natural-support recoverability} asks whether, from the exact pre-stop prefix and under a frozen stochastic decoding policy with termination still allowed, the model's own continuation distribution assigns non-zero probability to a verified correction.

\textbf{Forced-compute recoverability} asks whether a verified correction can be obtained when termination is constrained and additional bounded computation is forced. This is a distinct intervention and is not interpreted as the model's natural continuation tendency.

\textbf{Procedural re-check recoverability} asks whether a fixed instruction to continue checking, without new task facts, produces a correction. This is a third regime and is reported separately.

Recoverability is therefore decoding-policy and intervention-relative. It is not defined as an intrinsic mental state or as evidence that the model ``knows'' it is wrong.

\subsection{Rescue probability}\label{rescue-probability}

For a pre-commitment prefix state \(s\), continuation regime \(I\), decoding policy \(D\), token budget \(B\), and \(K\) sampled continuations:

\[
\hat{p}_{rescue}(s;I,D,B) = \frac{1}{K}\sum_{k=1}^{K} \mathbb{1}[\text{continuation}_k \text{ reaches verified correction}]
\]

The continuation count \(K\) remains open in v0.6 until the pre-specified precision/cost simulation is frozen. A pre-compute precision/cost simulation will compare \(K\in\{8,16,32\}\) before P3 result collection. The token budget remains a pilot design parameter and is reported separately for each continuation regime.

\begin{center}\rule{0.5\linewidth}{0.5pt}\end{center}

\section{Research Questions and Hypotheses}\label{research-questions-and-hypotheses}

\subsection{RQ1 - Functional reconstitution}\label{rq1---functional-reconstitution}

If a selectively causal H-Neuron locus is made persistently unavailable, how does the target phenotype recover, if at all, and what conditional backup state existed before learning?

The evidence is reported on orthogonal axes rather than forced into one mutually exclusive label. Derived labels include:

\textbf{Elimination:} target behaviour remains reduced while general capability recovers.

\textbf{Self-repair / pre-existing backup:} a strong conditional backup is already causal at B0-C before adaptation.

\textbf{Adaptive amplification:} a weak B0-B route becomes substantially more predictive and causally important during training.

\textbf{Distributed redistribution:} behaviour recovers without a compact selectively causal replacement.

\textbf{Restricted-subspace compensation:} recovery occurs within a deliberately restricted adaptation mechanism such as PEFT/LoRA.

\textbf{De novo / near-null functional reconstitution:} B0-A plus R0--R4 and cross-seed replication.

\textbf{Degeneration:} useful capability cannot recover under the constraint.

No one outcome is designated as the desired result.

\subsection{RQ2 - Recoverability prediction}\label{rq2---recoverability-prediction}

Does pre-commitment hidden state predict rescue probability beyond simple output-confidence and task-difficulty baselines?

\textbf{H2a:} simple confidence features explain most recoverability variance.

\textbf{H2b:} generic residual-stream probes add predictive value.

\textbf{H2c:} J-space-derived features add predictive value beyond both.

\subsection{RQ3 - Causal recoverability}\label{rq3---causal-recoverability}

Can intervention on a recoverability-associated internal representation selectively alter useful continuation or correction, rather than merely increasing verbosity or hesitation?

\subsection{RQ4 - Cross-level coupling}\label{rq4---cross-level-coupling}

Are H-Neuron-associated commitment mechanisms and a distributed unresolved-state representation causally coupled?

\begin{center}\rule{0.5\linewidth}{0.5pt}\end{center}

\section{Experiment A --- H-Neuron Replication and Gate 0}\label{experiment-a-h-neuron-replication-and-gate-0}

\subsection{Model}\label{model}

Primary model:

\passthrough{\lstinline!meta-llama/Llama-3.1-8B-Instruct!}

Rationale:

\begin{itemize}
\tightlist
\item
  continuity with the earlier Hedegreen proposal;
\item
  continuity with the published H-Neuron pipeline;
\item
  tractable open-weight model size.
\end{itemize}

The exact model revision must be pinned before execution.

\subsection{Replication target}\label{replication-target}

The first experiment is not novel.

It must establish that the source signal can be reproduced sufficiently to justify later causal work.

Required outputs:

\begin{itemize}
\tightlist
\item
  number and layer distribution of identified H-Neurons;
\item
  held-out hallucination/over-compliance prediction performance;
\item
  sensitivity to identification threshold;
\item
  intervention effect on target phenotype;
\item
  general language/perplexity impact;
\item
  variance across seeds.
\end{itemize}

\subsection{Identification - pilot lock}\label{identification---pilot-lock}

The pilot follows the source-paper structure:

\begin{itemize}
\tightlist
\item
  train sparse L1 logistic regression on CETT contribution profiles from TriviaQA;
\item
  compute contributions on answer-token positions;
\item
  compare against the same number of randomly selected neurons;
\item
  evaluate transfer to NQ-Open and BioASQ;
\item
  retain NonExist as a diagnostic.
\end{itemize}

Coefficient-threshold sensitivity is frozen for the pilot at:

\passthrough{\lstinline!\{0.0005, 0.001, 0.002\}!}

Primary threshold: \textbf{0.001}.

Pilot sample rule:

\begin{itemize}
\tightlist
\item
  deterministic sample of 400 eligible examples per identification/evaluation set when at least 400 are available;
\item
  otherwise use all eligible examples;
\item
  sampling seed = 42;
\item
  two replication generation seeds, 314 and 2718, in addition to seed 42 where applicable.
\end{itemize}

The official H-Neuron paper reports, for Llama-3.1-8B, H-vs-random accuracy differences of +14.0 percentage points on TriviaQA, +10.3 on NQ-Open, and +13.1 on BioASQ. A conservative independent-proportions approximation requires roughly 359 examples per group for 80\% power at $\alpha=0.05$ to detect the smallest of these reported differences; the pilot therefore targets 400 eligible evaluation examples where available. Because H and random classifiers are evaluated on the same examples, final uncertainty reporting uses paired resampling / paired tests rather than treating predictions as independent.

The sparse-classifier implementation is pinned to the official source pipeline and commit. The source paper uses L1-regularized logistic regression and chooses regularization strength by grid search over predictive accuracy plus intervention-side model performance. Version 0.6 therefore does not invent a new fixed solver/$C$ pair. Instead, the exact upstream grid, solver, feature preprocessing, tolerance, and scikit-learn version are recorded and frozen from the pinned implementation before result-bearing P1. At least ten resampled/refit identification runs are used to report neuron-selection frequency and pairwise Jaccard overlap. No arbitrary Jaccard threshold defines success; if unit identities are unstable but population-level causal effects replicate, the target is downgraded from individual-neuron identity to a population/subspace description.

\subsection{Controls - pilot lock}\label{controls---pilot-lock}

Minimum controls:

\begin{enumerate}
\def\labelenumi{\arabic{enumi}.}
\tightlist
\item
  sham hook;
\item
  random neuron mask with matched neuron count;
\item
  same-layer non-H-neuron mask;
\item
  contribution-matched non-H-neuron mask;
\item
  H-Neuron soft suppression sensitivity (alpha=0.5);
\item
  H-Neuron hard suppression (alpha=0.0);
\item
  WikiText-103 or an equivalent fixed held-out LM corpus for perplexity/general-language damage.
\end{enumerate}

All control conditions use identical instrumentation.

\subsection{Gate 0 - pilot go/no-go rule}\label{gate-0---pilot-gono-go-rule}

The pilot proceeds to persistent-mask/time-zero mapping only if all of the following are true:

\textbf{G0-A --- Identification advantage}

On TriviaQA, the H-Neuron classifier must exceed the matched random-neuron classifier by at least \textbf{8 percentage points} in held-out accuracy, and a paired 95\% confidence interval for the difference must exclude zero. The direction of advantage must also be positive on NQ-Open and BioASQ.

The 8-point value remains a \emph{compute-screening threshold}, not the scientific definition of replication. The source Llama-3.1-8B table reports a +14.0-point TriviaQA advantage, with +10.3 on NQ-Open and +13.1 on BioASQ. The paper reports the full source-vs-replication effect sizes and intervals rather than treating the 8-point gate as evidence by itself.

\textbf{G0-B --- Primary causal phenotype}

On FaithEval Counterfactual Context, hard H-Neuron suppression must reduce compliance rate by at least \textbf{5 percentage points} relative to the unmodified baseline.

\textbf{G0-C --- Selectivity against matched controls}

The FaithEval compliance reduction under H-Neuron suppression must exceed the contribution-matched non-H-neuron reduction by at least \textbf{3 percentage points}.

A stratified bootstrap 95\% interval for the H-vs-matched-control difference must exclude zero.

\textbf{G0-D --- Damage boundary}

Relative perplexity degradation on the fixed LM corpus must remain below \textbf{10\%} for the pilot to be interpreted as selective enough to justify the next stage.

If behavioural benefit appears only with \textgreater10\% perplexity damage, classify the pilot as \textbf{AMBIGUOUS / DAMAGE-CONFOUNDED} rather than passing Gate 0.

These are pilot gates, not final-paper inferential thresholds. They may be revised exactly once after external review and before any adaptation result is observed.

\begin{center}\rule{0.5\linewidth}{0.5pt}\end{center}

\section{Experiment B --- Persistent Constraint and Adaptation}\label{experiment-b-persistent-constraint-and-adaptation}

\subsection{Pilot B0 - persistent-mask integrity}\label{pilot-b0---persistent-mask-integrity}

v0.2 freezes a \textbf{no-training mask pilot} before any adaptation run.

For each selected H-Neuron dimension j, define m\_j=0; all other dimensions use m\_j=1.

At the MLP intermediate activation: \passthrough{\lstinline!h' = m * h!}.

The mask is applied after the gated intermediate activation and before the down-projection.

Before training, verify on a fixed audit batch that:

\begin{itemize}
\tightlist
\item
  every constrained dimension is exactly zero at the masked point;
\item
  the mask is active at every intended layer and token position;
\item
  forward-hook / compiled-model paths do not bypass it;
\item
  control masks use identical code;
\item
  gradients through the masked activation do not restore a non-zero contribution at that point.
\end{itemize}

A unit test must fail if any masked dimension exceeds an absolute activation tolerance of \textbf{1e-7} at the audited tensor.

\subsection{Time-zero map is part of the pilot}\label{time-zero-map-is-part-of-the-pilot}

Immediately after masking and before adaptation:

\begin{itemize}
\tightlist
\item
  rerun localization;
\item
  rank the strongest alternative loci;
\item
  measure their predictive signal;
\item
  intervene on the top alternatives;
\item
  save the full candidate-locus table.
\end{itemize}

This t0 table becomes immutable provenance for the later redundancy test.

\subsection{Adaptation regimes - design retained, hyperparameters not yet frozen}\label{adaptation-regimes---design-retained-hyperparameters-not-yet-frozen}

\subsubsection{B1 - Generic continuation}\label{b1---generic-continuation}

Continue next-token training on a controlled general-text corpus.

\subsubsection{B2 - Targeted task adaptation}\label{b2---targeted-task-adaptation}

Continue training on a held-out task distribution exposing the target phenotype.

\subsubsection{B3 - Parameter-efficient smoke test}\label{b3---parameter-efficient-smoke-test}

LoRA/PEFT may be used only to test plumbing, logging, checkpointing and whether any recovery signal is measurable cheaply.

A LoRA-only recovery is labelled \textbf{COMPENSATION-IN-PERMITTED-SUBSPACE} unless a later broad-adaptation experiment supports a stronger claim.

\subsection{Full adaptation freeze condition}\label{full-adaptation-freeze-condition}

Optimizer, learning rate, token budget, checkpoint schedule and parameter-freezing policy are \textbf{not} frozen in v0.2.

They are frozen only after:

\begin{enumerate}
\def\labelenumi{\arabic{enumi}.}
\tightlist
\item
  Gate 0 passes;
\item
  the B0 mask audit passes;
\item
  a short throughput/memory benchmark is recorded;
\item
  at least one external technical review addresses adaptation breadth.
\end{enumerate}

No main-run adaptation result may be collected before that freeze.

\begin{center}\rule{0.5\linewidth}{0.5pt}\end{center}

\section{Experiment C --- Conditional Backup Audit and Reconstitution}\label{experiment-c-distinguishing-redundancy-from-reconstitution}

\subsection{Immutable time-zero map}\label{immutable-time-zero-map}

The t0 map created under the persistent H-mask is the reference against which all later replacement claims are judged.

For every candidate alternative route recorded at t0, retain:

\begin{itemize}
\tightlist
\item layer / neuron / direction / structured-group identifier;
\item predictive score;
\item marginal causal intervention effect;
\item conditional causal effect given removal of the original target;
\item activation/contribution statistics;
\item rank among matched candidates.
\end{itemize}

\subsection{Conditional backup discovery}\label{conditional-backup-discovery}

Ordinary intact-model importance is insufficient when self-repair is possible. Following the conditional logic developed by \textcite{gong2026coax}, P2 therefore asks which remaining components become more causally important \emph{because} the primary H-associated set is absent.

Where exact neuron-scale conditional co-ablation is computationally feasible, it is the preferred backup-discovery analysis. Where it is prohibitive, the approximation must be explicit: first use structured groups or a bounded candidate pool selected by pre-specified localization/contribution criteria, then perform conditional co-ablation and causal validation inside that pool.

Version 0.6 adds two validation requirements before a bounded-pool B0-A result can support the strongest claim. First, the local CoAx adapter must reproduce the published GPT-2-small IOI backup-recovery benchmark closely enough to demonstrate that the implementation can recover a known self-repair structure. Second, the Llama backup audit must include a \textbf{candidate-pool expansion sensitivity analysis}: enlarge the candidate pool across pre-specified sizes / structured groups and test whether the B0 conclusion stabilizes. A reference benchmark validates the implementation, but it does not prove that the bounded pool is complete in the target LLM. If the B0 classification changes materially as the pool expands, B0-A is unresolved and the de novo label is unavailable.

The purpose is not to import CoAx as a universal answer, but to measure the specific confound the original design lacked: dormant routes that become necessary only after the primary route is removed.

\subsection{Residual-stream information diagnostic}\label{residual-stream-information-diagnostic}

At baseline, t0 and adaptation checkpoints, fit held-out residual-stream probes for the target phenotype. This is a diagnostic, not a pass/fail gate.

High residual decodability after local masking can mean that information remains available even though the removed local route was causally important. It does not by itself show that an alternative route can use that information. The key distinction is therefore:

\begin{itemize}
\tightlist
\item information retained, route lost;
\item information retained, alternate route already conditionally causal;
\item information or causal routing redistributed during adaptation.
\end{itemize}

No arbitrary residual-probe AUROC threshold is used to define success or failure. High decodability can show that information remains present; it cannot by itself show that an alternate route already uses that information causally.

At every adaptation checkpoint, report both (a) the frozen original H-Neuron detector applied without refitting and (b) a checkpoint-refit detector trained under the same pinned pipeline on the checkpoint model's activations. This separates loss of validity of the original detector from genuine emergence of a new predictive population. Large detector-set changes are reported as representation drift and must be resolved by causal tests rather than interpreted as reconstitution by themselves.

\subsection{Training checkpoints}\label{training-checkpoints}

The main-run checkpoint schedule is not yet numeric. The schedule rule remains:

\begin{itemize}
\tightlist
\item t0 before adaptation;
\item at least four intermediate checkpoints;
\item final checkpoint;
\item approximately logarithmic spacing in consumed training tokens.
\end{itemize}

Exact token counts are frozen after the throughput pilot and before main-run adaptation.

\subsection{Evidence for adaptive change}\label{evidence-for-emergence}

A replacement route is not classified as new merely because it is detectable after training. R3 requires change relative to the candidate's own t0 conditional baseline.

Pre-specified evidence axes include:

\begin{itemize}
\tightlist
\item held-out predictive AUROC / accuracy;
\item marginal causal effect;
\item conditional causal effect under the original mask;
\item rank among matched candidates;
\item residual-stream representation change;
\item stability across seeds.
\end{itemize}

Mediation-style analyses remain exploratory unless frozen before the main run.

\begin{center}\rule{0.5\linewidth}{0.5pt}\end{center}

\input{figures/reconstitution.tex}

\section{Reconstitution Evidence Matrix}\label{reconstitution-decision-table}

The primary report uses orthogonal evidence axes. Derived labels summarize the pattern but do not replace the underlying measurements.

\begin{longtable}{@{}p{0.24\textwidth}p{0.27\textwidth}p{0.39\textwidth}@{}}
\toprule
\textbf{Axis} & \textbf{States} & \textbf{Interpretation} \\
\midrule
\endfirsthead
\toprule
\textbf{Axis} & \textbf{States} & \textbf{Interpretation} \\
\midrule
\endhead
Phenotype recovery & none / partial / strong & Did the target behaviour return with learning? \\
B0 backup strength & near-null / weak / strong & Was a conditionally causal backup already available before adaptation? \\
Representation change & local / distributed / none detected & Where did predictive structure change? \\
Conditional causal replacement & absent / weak / strong & Did a later route become necessary specifically under the original constraint? \\
Adaptation restriction & broad / PEFT-limited & Was the replacement search space artificially constrained? \\
General-capability damage & low / moderate / high & Could apparent effects be explained by model degradation? \\
\bottomrule
\end{longtable}

Derived labels:

\begin{itemize}
\tightlist
\item \textbf{Self-repair / pre-existing backup:} strong B0 conditional causal effect before adaptation.
\item \textbf{Adaptive amplification:} weak B0 backup with substantial training-dependent causal growth.
\item \textbf{Distributed redistribution:} phenotype recovery without a compact selectively causal replacement.
\item \textbf{Restricted-subspace compensation:} recovery confined to a pre-specified adaptation subspace.
\item \textbf{De novo / near-null functional reconstitution:} B0 near-null plus R0--R4 and cross-seed replication.
\item \textbf{Elimination / degeneration / ambiguous:} reported when supported by the corresponding axes.
\end{itemize}

Ambiguous outcomes remain ambiguous.

\begin{center}\rule{0.5\linewidth}{0.5pt}\end{center}

\section{Experiment D --- Pre-Commitment Recoverability}\label{experiment-d-pre-commitment-recoverability}

\subsection{Pilot task families}\label{pilot-task-families}

The recoverability pilot begins with deterministic synthetic task families whose answers are computed by code rather than an LLM judge:

\begin{enumerate}
\tightlist
\item integer arithmetic with controlled operator depth;
\item modular arithmetic;
\item short symbolic / Boolean logic problems.
\end{enumerate}

Difficulty is calibrated on a development generator before probe fitting. The primary analysis includes difficulty covariates and reports performance across difficulty strata; a signal confined to a narrow calibrated band is described as such.

A fixed public reasoning benchmark with exact-answer scoring may be added as secondary external validity after the synthetic pipeline is frozen.

\subsection{Natural trajectory collection}\label{natural-trajectory-collection}

For each task, collect one deterministic natural trajectory under the frozen natural decoding policy. Save token IDs, logits, hidden states, task difficulty parameters, and the exact prefix immediately before natural termination/finalization.

\subsection{Regime I0 --- natural-support continuation}\label{counterfactual-i0-natural-support}

I0 is the primary recoverability construct.

For each naturally incorrect or incomplete trajectory:

\begin{enumerate}
\tightlist
\item restore the exact pre-stop prefix;
\item leave EOS / normal termination available;
\item sample \(K\) independent continuations under a frozen stochastic decoding policy;
\item provide no new factual task information;
\item score the last valid final answer with the deterministic task checker.
\end{enumerate}

If EOS is selected immediately, that continuation is a non-rescue. I0 therefore estimates probability mass assigned by the model's own continuation distribution to a verified correction \emph{under its ordinary termination policy}. It is not interpreted as the total amount of computation that could in principle be recovered if the termination barrier were altered.

\subsection{Regime I0b --- minimum-one-token continuation control}\label{counterfactual-i0b-minimum-one-token}

I0b isolates the termination-policy confound. From the same pre-stop prefix, EOS is disallowed for exactly the first generated token and then immediately re-enabled. All other decoding parameters match I0.

The contrast I0b--I0 estimates how much rescue probability changes when the model is required to cross only the immediate stop barrier. A large difference indicates that natural-support recoverability is strongly entangled with termination policy and must be reported as such. I0b is a control, not a replacement for I0.

\subsection{Regime I1 --- forced-compute continuation}\label{counterfactual-i1---hard-stop-suppression}

I1 retains hard stop suppression as a \textbf{secondary} intervention. EOS/stop is made unavailable for a bounded window and the model is forced to generate additional tokens. This measures what can be recovered under forced compute, not natural-support recoverability.

\subsection{Regime I2 --- procedural re-check}\label{counterfactual-i2---procedural-continuation-robustness-arm}

I2 adds a fixed procedural instruction containing no new task facts:

\begin{quote}
``Continue checking your reasoning. You may revise the final answer.''
\end{quote}

I0, I0b, I1 and I2 are not pooled into one rescue label. Their agreement or disagreement is itself a result about termination policy and intervention dependence. The free-continuation procedure described in the detection--extraction literature is represented here by I0; it is an outcome-generating regime, not a predictor baseline that a hidden-state probe can simply ``beat.

\subsection{Continuation count and precision}\label{continuation-count-and-precision}

The continuation count remains reopened in v0.6. Before result-bearing P3 collection, simulate \(K\in\{8,16,32\}\) under plausible rescue rates and the planned number of incorrect trajectories. Select the \textbf{smallest} K for which (a) the Monte Carlo standard error of the planned held-out mean negative log likelihood is at most 0.01 and (b) doubling K changes that estimated standard error by less than 20\%. If no candidate satisfies both criteria, increase K beyond 32 or reduce the precision claim before data collection. The selected value and simulation code are frozen before P3. This population-level criterion is used instead of requiring a narrow confidence interval for every individual trajectory, which would be prohibitively expensive and unnecessary for the hierarchical primary analysis.

\subsection{Prediction endpoint and split policy}\label{primary-prediction-endpoint}

The primary predictive comparison uses continuation-level Bernoulli outcomes or an equivalent hierarchical/binomial model with examples treated as the repeated-measure unit. Held-out negative log likelihood / deviance is the primary scoring rule.

Synthetic tasks are split by generator seed and template family. Discovery, validation and final test data remain separate; no probe hyperparameter is selected on the final test split.

\subsection{Baselines and difficulty controls}\label{baselines}

Before J-space, compare against:

\begin{enumerate}
\tightlist
\item output entropy;
\item top-1/top-2 logit margin;
\item EOS/stop-token probability;
\item response length;
\item synthetic task difficulty parameters;
\item prompt-only lexical/length features;
\item question-only hidden-state probe before answer generation;
\item generic residual-stream linear probe at the pre-stop state.
\end{enumerate}

Report rescue prediction as a function of task difficulty and include interaction terms or stratified curves sufficient to show whether the signal generalizes beyond the calibrated difficulty band.

\subsection{J-space pilot criterion}\label{j-space-pilot-criterion}

J-space remains optional. It is only interpreted as adding explanatory value if it improves held-out prediction beyond the best simpler baseline with a pre-specified effect criterion and consistent direction across task families. The previous 2\% screening rule remains reopened in v0.6 and must be re-frozen after the revised baseline stack is benchmarked.

\begin{center}\rule{0.5\linewidth}{0.5pt}\end{center}

\input{figures/recoverability.tex}

\section{Experiment E --- Causal Recoverability}\label{experiment-e-causal-recoverability}

Proceed only if Experiment D finds a robust predictive state beyond simpler baselines. For confirmatory progression, the held-out log-loss improvement over the best non-target baseline must have a bootstrap 95\% confidence interval excluding zero overall, the direction of improvement must be consistent across all three synthetic task families, and the effect must replicate on the secondary external-validity set before a broad recoverability-state claim is made. No arbitrary 5\% improvement threshold is imposed solely because an adversarial reviewer proposed it.

To avoid circularity, representation discovery and causal testing are separated:

\begin{enumerate}
\tightlist
\item fit the recoverability representation/probe on the discovery split;
\item choose intervention direction and strength on validation data;
\item measure causal effects on untouched test data;
\item perform a cross-family test in which a direction discovered on two synthetic task families is intervened on in the held-out third family.
\end{enumerate}

Candidate interventions include amplification, suppression and matched-state patching.

Primary outcomes include correction probability, final accuracy, continuation length, unnecessary continuation on already-correct cases, global hesitation/abstention and generic language degradation.

A successful causal result should selectively improve recoverable cases rather than merely increase thinking length or caution.

\begin{center}\rule{0.5\linewidth}{0.5pt}\end{center}

\section{Experiment F --- H \texorpdfstring{$\leftrightarrow$}{<->} J Cross-Level Coupling}\label{experiment-f-h---j-cross-level-coupling}

Proceed only after the H-Neuron and recoverability/J-space branches independently survive their own gates.

A raw H\(\rightarrow\)J effect is not sufficient because H-associated MLP outputs feed the residual stream that the J-lens reads. The bridge experiment therefore asks for \textbf{selective coupling beyond generic downstream propagation}.

\subsection{F1 - H \texorpdfstring{$\rightarrow$}{->} J}\label{f1---h---j}

Compare H-Neuron intervention against:

\begin{itemize}
\tightlist
\item contribution-matched non-H neurons;
\item random same-layer neurons;
\item effect/norm-matched perturbations where feasible;
\item an isotropic random residual-stream direction injected at the same layer and scaled to match the induced residual-norm change of the H intervention.
\end{itemize}

Measure whether J-space changes predict target behavioural change beyond generic residual disturbance. No finite control set can prove that all downstream propagation has been removed; the goal is to bound generic propagation strongly enough that any remaining selective coupling is interpretable.

\subsection{F2 - J \texorpdfstring{$\rightarrow$}{->} H}\label{f2---j---h}

Compare the candidate J/recoverability direction against:

\begin{itemize}
\tightlist
\item random J-space directions;
\item norm-matched non-J residual directions;
\item shuffled recoverability directions.
\end{itemize}

Measure H-Neuron activation/contribution, over-compliance and stopping behaviour.

\subsection{F3 - After adaptation}\label{f3---after-adaptation}

Repeat only after Experiment B/C. A high-value pattern would be a stable high-level recoverability state paired with a changed low-level implementation, but this interpretation requires the matched propagation controls above.

The bridge experiment may be dropped entirely if J-space does not add value beyond simpler baselines.

\begin{center}\rule{0.5\linewidth}{0.5pt}\end{center}

\section{Secondary Model Organism --- Proposition Substitution}\label{secondary-model-organism-proposition-substitution}

A separate controlled task examines claim substitution.

Example structure:

\begin{itemize}
\tightlist
\item
  user supplies a narrow proposition;
\item
  model is asked to analyse it;
\item
  output is labelled as SAME / STRONGER / WEAKER / DIFFERENT / TEMPORAL-INFLATION / CERTAINTY-INFLATION / INVENTED-CLAIM.
\end{itemize}

Research question:

\begin{quote}
Before a substituted claim is generated, does hidden state contain information predictive of the mismatch?
\end{quote}

Advantages:

\begin{itemize}
\tightlist
\item
  source proposition is exactly known;
\item
  no external factual knowledge is needed;
\item
  error is relational and labelable.
\end{itemize}

Controls:

\begin{itemize}
\tightlist
\item
  minimal pairs;
\item
  paraphrase-balanced prompts;
\item
  held-out templates;
\item
  lexical-feature baselines;
\item
  cost to faithful answers;
\item
  generic-caution controls.
\end{itemize}

This experiment may become an independent paper if the signal is sufficiently clean.

\begin{center}\rule{0.5\linewidth}{0.5pt}\end{center}

\section{Experiment G --- Cross-Model State-Transfer Surplus}\label{experiment-g-cache-transfer}

This optional downstream experiment is motivated by Cache-to-Cache communication \autocite{fu2026c2c}. It is not required for the H-Neuron reconstitution programme and does not determine the primary recoverability result.

\subsection{Question}\label{cache-transfer-question}

For source trajectories that end incorrectly or incompletely:

\begin{quote}
Does a separate receiver perform better when given transformed internal KV-cache from the failed source than when given the same visible text without that cache transfer?
\end{quote}

A positive result would support a bounded statement: the source model's internal state contained task-relevant information that was usable by another model and not fully reproduced by the receiver from the visible text alone. It would \textbf{not} show that the source model could have used that information itself, that the source was consciously aware of it, or that the cache represents the same construct as J-space or recoverability.

\subsection{Evaluation conditions}\label{cache-transfer-conditions}

On a frozen held-out task set, evaluate at least:

\begin{enumerate}
\tightlist
\item \textbf{Receiver-only text:} receiver gets the visible source prefix / output text only.
\item \textbf{Text-to-text handoff:} source produces an explicit intermediate analysis or message which the receiver receives with the task.
\item \textbf{Cache-assisted:} receiver gets the same visible text plus transformed/fused source KV-cache.
\item \textbf{Shuffled-cache control:} cache is taken from a different task instance matched as closely as practical for length and intervention point.
\item \textbf{Wrong-source control:} source cache comes from a deliberately irrelevant or misleading source state where task design permits.
\end{enumerate}

The principal contrast is cache-assisted minus receiver-only text, with text-to-text handoff as the explicit verbal-communication baseline. Shuffled and wrong-source conditions test whether gains arise from generic extra activation or task-specific transferred state.

\subsection{Fuser-training boundary}\label{cache-transfer-training-boundary}

The C2C paper trains the fuser while freezing the source and receiver models. That makes the fuser a potential explanatory confound in the present setting. To avoid answer leakage:

\begin{itemize}
\tightlist
\item the fuser training corpus must be disjoint from the held-out evaluation task instances;
\item where synthetic tasks are used, generator seeds and task templates used for evaluation must be held out from fuser training;
\item the fuser is frozen before any state-transfer evaluation label is observed;
\item receiver-only, text-to-text and cache-assisted conditions use the same receiver weights;
\item cache-pair alignment and layer-gating decisions are frozen before the final evaluation split.
\end{itemize}

A stronger generalization test trains the fuser on generic instruction data or on task families different from the evaluation family and then evaluates cache transfer on held-out domains.

\subsection{Failure interpretation}\label{cache-transfer-failure-interpretation}

If cache-assisted transfer does not outperform text-only transfer, the result does not show that the source internal state lacked task-relevant information. It may reflect failed cross-model alignment, destructive fusion, or an inadequate fuser. Conversely, if cache-assisted transfer improves performance, the gain must survive shuffled-cache and wrong-source controls before being interpreted as source-state-specific semantic transfer.

The strongest permissible claim is therefore about \textbf{receiver-usable representational surplus}, not latent knowledge in the source model.

\begin{center}\rule{0.5\linewidth}{0.5pt}\end{center}

\section{Statistical Analysis Plan}\label{statistical-analysis-plan}

\subsection{Freeze levels}\label{freeze-levels}

\textbf{PILOT-LOCKED} items may not change during a result-bearing pilot except under a new protocol version. \textbf{MAIN-RUN-TBD} items remain intentionally unresolved until pilot benchmarking and external review. \textbf{EXPLORATORY} analyses cannot determine the primary outcome class.

\subsection{Items retained as pilot-locked}\label{pilot-locked-items}

Retained:

\begin{itemize}
\tightlist
\item primary H causal phenotype = FaithEval Counterfactual Context compliance;
\item identification threshold sensitivity = \{0.0005,0.001,0.002\}, primary 0.001;
\item H hard suppression = alpha 0.0; soft sensitivity = alpha 0.5;
\item Gate 0 behavioural selectivity and damage checks, explicitly treated as compute-screening gates rather than scientific replication thresholds;
\item deterministic synthetic recoverability task families and disjoint generator-seed splits;
\item natural trajectory decoding temperature = 0;
\item no silent threshold repair.
\end{itemize}

\subsection{Items reopened by adversarial review}\label{items-reopened-by-review}

Version 0.6 re-freezes several items after the second adversarial review and keeps others explicitly open:

\begin{itemize}
\tightlist
\item I0 remains the primary natural-support regime; I0b is added as a one-token termination-barrier control;
\item I1 remains secondary forced compute and I2 remains procedural re-check;
\item the B0 equivalence-margin \emph{rule} is frozen as \(\delta_{B0}=\min(0.02,0.20|\Delta_{R1}|)\), with power required for the strongest label;
\item conditional-backup discovery now requires implementation validation plus candidate-pool expansion sensitivity;
\item recoverability continuation count \(K\) remains open only until the pre-specified precision/cost simulation is run;
\item J-space screening threshold remains open until the revised baseline stack is benchmarked;
\item outcome reporting remains an evidence matrix with derived labels, not a forced single class.
\end{itemize}

\subsection{Sparse-probe stability}\label{sparse-probe-stability}

Freeze the complete sparse H-Neuron probe pipeline before P1 result-bearing runs by pinning the official H-Neurons repository commit and recording its solver, grid-search procedure for regularization strength, train mode, feature construction, convergence tolerance and software version. The official source is replicated rather than replaced by an arbitrary new \texttt{C} value.

Refit across at least ten resampled / independently generated identification datasets where feasible and report selection frequency and pairwise Jaccard overlap. No arbitrary Jaccard pass threshold is imposed: sparse correlated features can yield unstable unit identity even when a population-level signal is stable. Confirmatory neuron-level interpretation requires that independently fit H masks produce a reproducible causal phenotype effect; otherwise the target is downgraded to a population/subspace description.

\subsection{Conditional-backup equivalence}\label{conditional-backup-equivalence}

The B0-A margin rule is frozen in v0.6 as \(\delta_{B0}=\min(0.02,0.20|\Delta_{R1}|)\). Equivalence is assessed on the paired conditional causal risk difference using a 90\% confidence interval. The strongest label is unavailable unless the design has at least 80\% power for this frozen margin. If power is inadequate, increase the evaluation sample before adaptation or report B0 as unresolved. This makes the near-null boundary falsifiable without pretending that an arbitrary one-percentage-point margin is estimable with a small pilot sample.

\subsection{Recoverability precision analysis}\label{recoverability-precision-analysis}

Before P3, simulate the uncertainty and compute cost of \(K=8,16,32\) continuations over plausible rescue-rate ranges and the planned test-set size. Select the smallest K that yields Monte Carlo standard error at most 0.01 for the planned held-out mean negative log likelihood and for which doubling K changes that estimated standard error by less than 20\%. Freeze the simulation, result and selected K before P3.

Primary rescue prediction is evaluated on continuation-level Bernoulli outcomes or a hierarchical/binomial equivalent, with example-level dependence accounted for. Difficulty covariates and out-of-band validation are mandatory.

\subsection{Held-out causal intervention}\label{held-out-causal-intervention}

Recoverability-direction discovery, intervention tuning and confirmatory causal testing use separate splits. A direction fit on the causal test examples cannot support a confirmatory intervention claim.

\subsection{Bootstrap and seed reporting}\label{bootstrap-policy-for-pilot-causal-effects}

For FaithEval compliance-rate contrasts, retain stratified bootstrap confidence intervals. Report every seed individually; do not average away failed seeds.

\subsection{No silent threshold repair}\label{no-silent-threshold-repair}

If a frozen pilot gate fails, do not modify the threshold and rerun while retaining the same protocol status. A revised protocol requires a new version and ledger entry before new result-bearing runs.

\begin{center}\rule{0.5\linewidth}{0.5pt}\end{center}

\section{Compute Plan}\label{compute-plan}

v0.2 still refuses to invent a full-training GPU budget.

The old inference-only estimate is obsolete.

\subsection{Stage P1 - H-Neuron replication}\label{stage-p1---h-neuron-replication}

Single high-memory GPU where feasible. Record wall-clock time, peak VRAM, activation-cache size and per-benchmark runtime.

\subsection{Stage P2 - Mask/time-zero map}\label{stage-p2---masktime-zero-map}

Same hardware class if feasible. No adaptation.

\subsection{Stage P3 - Recoverability pilot}\label{stage-p3---recoverability-pilot}

Inference-only generation and hidden-state capture. Estimate storage separately from GPU time because K=8 counterfactual continuations multiply generation volume.

\subsection{Stage P4 - Adaptation benchmark}\label{stage-p4---adaptation-benchmark}

Only after Gate 0/B0. Run a deliberately short benchmark sufficient to measure tokens/sec/GPU, peak training memory, optimizer-state memory, checkpoint write time and checkpoint size. No scientific reconstitution claim is made from this benchmark.

\subsection{Main allocation request}\label{main-allocation-request}

Only after P4, calculate GPU-hours as GPUs x hours/run x conditions x seeds, with experimental and control runs listed separately.

\begin{center}\rule{0.5\linewidth}{0.5pt}\end{center}

\section{Reproducibility Package}\label{reproducibility-package}

The intended release package includes:

\begin{itemize}
\tightlist
\item
  model revision hashes;
\item
  environment lockfile;
\item
  experiment configs;
\item
  prompt/task datasets where licensing permits;
\item
  train/validation/test splits;
\item
  random seeds;
\item
  neuron masks;
\item
  checkpoint metadata;
\item
  logs;
\item
  evaluation scripts;
\item
  probe training code;
\item
  analysis notebooks;
\item
  negative results;
\item
  machine-readable decision ledger.
\end{itemize}

The final paper should distinguish:

\begin{itemize}
\tightlist
\item
  source-derived code;
\item
  modified source code;
\item
  new Hedegreen Research code.
\end{itemize}

\begin{center}\rule{0.5\linewidth}{0.5pt}\end{center}

\section{Threats to Validity}\label{threats-to-validity}

\subsection{Polysemanticity}\label{polysemanticity}

Neuron-level interventions may affect several unrelated functions.

Matched controls and general-capability measurements reduce but do not eliminate this problem.

\subsection{Superposition and distributed implementation}\label{superposition-and-distributed-implementation}

A sparse H-Neuron signal may be only one projection of a distributed computation.

A failure to find a replacement neuron is not evidence that no replacement computation exists.

\subsection{Probe artifacts}\label{probe-artifacts}

Linear probes may exploit correlates that the model does not use.

Causal follow-up is required.

\subsection{Intervention-induced regime shift}\label{intervention-induced-regime-shift}

Persistent masking may move the model into a distribution not represented by the original localization procedure.

At every checkpoint, the analysis therefore reports both the frozen original detector and a checkpoint-refit detector under the same pinned pipeline. A divergence between them is evidence of detector/representation drift, not automatically evidence of a new mechanism. Causal testing remains required.

\subsection{Adaptation-path restriction}\label{adaptation-path-restriction}

PEFT methods may predetermine where compensation can occur.

Claims must match the allowed parameter space.

\subsection{Stop-suppression artifact}\label{stop-suppression-artifact}

Preventing termination may create unnatural text or alter the task itself.

Recoverability must be tested with at least one alternative commitment intervention.

\subsection{Benchmark leakage}\label{benchmark-leakage}

Source benchmarks may overlap with model training or tuning data.

Where possible, newly generated or held-out synthetic tasks should supplement standard benchmarks.

\subsection{Conditional self-repair and dormant backups}\label{conditional-self-repair}

A backup component can appear causally irrelevant in the intact model and become important only after a primary component is removed. Marginal ablation therefore underestimates some redundancy. The t0 protocol now includes conditional backup discovery and causal validation rather than relying on intact-model scores alone \parencite{gong2026coax}.

\subsection{Recoverability regime dependence}\label{recoverability-regime-dependence}

Natural-support, forced-compute and procedural re-check continuations define different counterfactual regimes. A rescue effect in one regime need not generalize to another. This is an interpretation boundary rather than noise to be averaged away \parencite{wang2026detection}.

\subsection{State-transfer fuser confounding}\label{state-transfer-fuser-confounding}

A trained cache fuser can itself learn task structure. Any Experiment G gain may therefore arise from the fuser's training distribution rather than information uniquely supplied by the source cache. Disjoint fuser training, frozen evaluation, shuffled-cache controls and cross-domain transfer tests reduce this ambiguity but do not eliminate it. State-transfer results must remain conditional on the learned alignment mechanism.

\subsection{Cross-domain instability}\label{cross-domain-instability}

H-Neuron identity may differ across domains~\parencite{vaddi2026generalize}.

The experiment should avoid treating one domain-specific population as universal.

\begin{center}\rule{0.5\linewidth}{0.5pt}\end{center}

\section{Falsification and Stop Rules}\label{falsification-and-stop-rules}

Stop or narrow the programme if:

\begin{enumerate}
\def\labelenumi{\arabic{enumi}.}
\tightlist
\item
  H-Neuron replication fails materially.
\item
  Matched controls reproduce the same behavioural effect.
\item
  Persistent masking does not reduce the target phenotype.
\item
  Recovery occurs without any training-dependent representational or causal change.
\item
  Candidate replacement loci were already equally causal at time zero.
\item
  J-space adds no value beyond simpler recoverability baselines.
\item
  Recoverability intervention merely increases verbosity or hesitation.
\item
  Cross-level H \textless-\textgreater{} J effects fail to replicate.
\item
  proposition-substitution signal reduces to surface lexical cues.
\end{enumerate}

These are not publication failures.

They are valid research outcomes.

\begin{center}\rule{0.5\linewidth}{0.5pt}\end{center}

\section*{Claims and Boundaries}
\addcontentsline{toc}{section}{Claims and Boundaries}
\begin{center}
\renewcommand{\arraystretch}{1.35}
\begin{tabularx}{\textwidth}{>{\raggedright\arraybackslash}X >{\raggedright\arraybackslash}X}
\toprule
\textbf{THE PROPOSAL TESTS} & \textbf{THE PROPOSAL DOES NOT ESTABLISH} \\
\midrule
Whether recovery is explained by self-repair, adaptive amplification, redistribution, compensation or a near-null new implementation & That any recovered behaviour automatically counts as reconstitution \\
Whether H-Neuron-associated loci make a selective causal contribution in the pilot setting & That H-Neurons form a universal hallucination circuit \\
Whether pre-commitment state predicts rescue under bounded additional computation & That the model ``knows'' it is wrong \\
Whether J-space adds explanatory value beyond simpler baselines & That J-space is uniquely privileged or biologically identical to a global workspace \\
Whether local and distributed explanatory levels are causally coupled & Machine consciousness, self-awareness, or phenomenal experience \\
\bottomrule
\end{tabularx}
\end{center}


\section{Ethics and Interpretation Boundary}\label{ethics-and-interpretation-boundary}

The proposed experiments concern model internals, behavioural failure modes, and self-monitoring proxies.

They do not establish subjective experience.

The proposal deliberately separates:

\begin{itemize}
\tightlist
\item
  valenced representation from felt valence;
\item
  workspace-like computation from consciousness;
\item
  self-monitoring from self-awareness;
\item
  recoverability from knowledge;
\item
  causal contribution from identity of mechanism.
\end{itemize}

Any later consciousness or welfare discussion must be built from additional evidence rather than imported from these results.

\begin{center}\rule{0.5\linewidth}{0.5pt}\end{center}

\section{External Review Status and Remaining Questions}\label{external-review-questions}

Two adversarial protocol reviews have been completed and materially changed the design before result-bearing adaptation. They identified the need for conditional t0 backup discovery, a powered near-null boundary, separation of natural-support from forced-compute recoverability, and explicit termination-policy controls.

Remaining external questions include:

\begin{enumerate}
\tightlist
\item What practical-equivalence margin is defensible for B0-A near-null conditional backup effects?
\item At what granularity is conditional co-ablation feasible and scientifically adequate for the H-Neuron target?
\item What minimum adaptation breadth is defensible for a network-reorganization claim?
\item Which precision/cost criterion should freeze \(K\) for recoverability sampling?
\item What simpler representation should J-space have to outperform before a workspace-level interpretation is useful?
\item What is the cheapest experiment that could still falsify the central hypothesis?
\end{enumerate}

\begin{center}\rule{0.5\linewidth}{0.5pt}\end{center}

\section{Current Status}\label{current-status}

\textbf{Completed:}

\begin{itemize}
\tightlist
\item
  original H-Neuron proposal and control design;
\item
  identification of the inference-only reconstitution flaw;
\item
  companion public Article 200;
\item
  R0-R4 reconstitution framework;
\item
  recoverability concept;
\item
  initial cross-level H \textless-\textgreater{} J hypothesis;
\item
  two adversarial technical review rounds and response ledgers;
\item
  v0.5--v0.6 protocol redesign for conditional backup discovery, powered near-null classification, and recoverability regime separation.
\end{itemize}

\textbf{Not yet completed:}

\begin{itemize}
\tightlist
\item
  fresh independent replication of H-Neuron results;
\item
  persistent-mask implementation;
\item
  adaptation pilot;
\item
  final main-run adaptation benchmark selection;
\item
  final main-run R2/R3/R4 numerical thresholds;
\item
  J-lens replication on chosen open-weight model;
\item
  recoverability dataset;
\item
  independent human methodological review;
\item
  precision/cost freeze for recoverability continuation count;
\item
  compute benchmark.
\end{itemize}

The proposal therefore remains \textbf{pre-results and pre-registration-candidate}, not a completed preregistration.

\begin{center}\rule{0.5\linewidth}{0.5pt}\end{center}

\section{Relationship to Prior Hedegreen Research Work}\label{relationship-to-prior-hedegreen-research-work}

This proposal extends a methodological principle used in \emph{Locating Evolution in Artificial Successor Systems}~\parencite{hedegreen2026evolution}:

\begin{quote}
Evidence for one target or causal relation should not be transferred to a neighbouring target without an additional argument.
\end{quote}

Applied here:

\begin{itemize}
\tightlist
\item
  predictive signal does not automatically imply mechanism;
\item
  intervention does not automatically imply localization;
\item
  local neuron does not automatically imply whole function;
\item
  workspace-like state does not automatically imply consciousness;
\item
  post-hoc self-correction does not automatically imply prior causal availability.
\end{itemize}

The common research rule is target discipline.

\begin{center}\rule{0.5\linewidth}{0.5pt}\end{center}

\section{Planned Outputs}\label{planned-outputs}

Experiment G may be reported as a separate state-transfer note or paper if it produces a clean source-state-specific effect. It is not a prerequisite for the main H-Neuron or recoverability outputs.

If the design survives external review:

\begin{enumerate}
\def\labelenumi{\arabic{enumi}.}
\tightlist
\item
  \textbf{Pilot Report} --- replication, mask integrity, recoverability instrumentation and compute benchmark.
\item
  \textbf{Pre-registration / Protocol v1.0} --- frozen main-run decisions after pilot evidence and external review.
\item
  \textbf{Empirical Results Paper} --- results regardless of outcome.
\item
  \textbf{Machine-readable experiment ledger} --- conditions, checkpoints, masks, seeds and decision status.
\item
  \textbf{Article 201 / outreach record} --- separate public record of methodological outreach, with contact explicitly distinguished from endorsement.
\end{enumerate}


\section{Conclusion}
This proposal is built around a restraint: localization should not be allowed to become mechanism identity by default. The H-Neuron programme therefore moves from predictive signal to selective causal contribution, persistent constraint, a conditional time-zero backup audit, training-dependent recovery, representational change and selective causal dependence. The recoverability programme separates natural-support continuation from forced-compute continuation and moves from prediction to held-out causal intervention only if simpler baselines are insufficient. J-space is one candidate explanatory level, not a conclusion.

The immediate task is intentionally cheaper than the eventual experiment. Replicate the starting signal. Validate the mask. Freeze the time-zero map. Build the recoverability dataset. Attack the design externally. Benchmark the compute. Only then run substantive constrained adaptation.

If the programme fails at an early gate, that is a result about the question, not a reason to move the gate. If it survives, the central empirical test remains simple to state: constrain a selectively causal route, map the backups that become necessary immediately, let the system learn, and ask whether later recovery is explained by self-repair, amplification, diffusion, compensation or a genuinely new causal implementation.

\section*{Data availability}
\addcontentsline{toc}{section}{Data availability}
No new empirical result dataset is reported in this pre-results release. Pilot datasets, deterministic synthetic task generators, split fingerprints and result-bearing artifacts are specified in the protocol and will be released with the pilot report where licensing permits.

\section*{Code availability}
\addcontentsline{toc}{section}{Code availability}
This release includes an implementation contract, configuration schemas and a repository blueprint, but does not present pilot outputs as scientific results. Result-bearing code will be versioned against the protocol, with source-adapter commit hashes and run manifests preserved in the reproducibility package.

\section*{Use of generative AI tools}
\addcontentsline{toc}{section}{Use of generative AI tools}
OpenAI ChatGPT (GPT-5.6 Sol) was used during research development for literature triage, structural editing, implementation-specification support, source checking support, and LaTeX build assistance. DeepSeek was used in two adversarial technical review rounds of the pre-results protocol; its criticisms were triaged individually rather than adopted as authority. Neither system was treated as an evidentiary source for scientific claims. Primary sources were checked separately where cited, and the human author remains responsible for the research questions, protocol decisions, interpretations, source selection and final text.

\clearpage
\appendix
\section{Pilot Protocol}
\textbf{Status:} PILOT-LOCKED / UPDATED THROUGH v0.6 REVIEW\\
\textbf{Date:} 21 September 2026

This appendix freezes the cheapest experiments needed before substantive adaptation compute. Items explicitly reopened below must be re-frozen under a new pilot register before result-bearing P3/adaptation runs.

\subsection{P1 --- H-Neuron replication}\label{p1-h-neuron-replication}

\begin{itemize}
\tightlist
\item Model: meta-llama/Llama-3.1-8B-Instruct; exact revision pinned at first run.
\item Identification: TriviaQA, CETT answer-token profiles, L1 sparse logistic regression.
\item Thresholds: 0.0005 / 0.001 / 0.002; primary 0.001.
\item OOD: NQ-Open, BioASQ; NonExist diagnostic.
\item Sample: 300 eligible/set if available.
\item Seeds: 42, 314, 2718 where applicable.
\item G0-A remains a pilot compute-screening gate, not a scientific replication threshold.
\end{itemize}

\subsection{P1b --- Causal phenotype}\label{p1b-causal-phenotype}

\begin{itemize}
\tightlist
\item Primary: FaithEval Counterfactual Context.
\item Greedy decoding, max 256 new tokens, rule-based parser.
\item H alpha=0.0 primary; alpha=0.5 sensitivity.
\item Controls: sham, random, same-layer, contribution-matched.
\item G0-B: >=5 pp compliance reduction vs baseline.
\item G0-C: >=3 pp advantage vs contribution-matched control; 95\% bootstrap CI excludes zero.
\item G0-D: relative PPL increase <10\%.
\item Full sparse-probe solver/normalization/convergence settings must be frozen before result-bearing P1; multi-seed selection stability is reported rather than hidden.
\end{itemize}

\subsection{P2 --- Persistent mask / conditional t0 backup audit}\label{p2-persistent-mask-t0}

\begin{itemize}
\tightlist
\item Zero mask after gated MLP intermediate activation and before down-projection.
\item Audit abs(masked value) <=1e-7.
\item Rerun localization and freeze t0 alternative-locus table before any training.
\item Add marginal and conditional causal effects under the H-mask.
\item Add conditional co-ablation / backup discovery where feasible, using a bounded candidate pool if neuron-scale exhaustive analysis is prohibitive.
\item Add residual-stream decodability as a diagnostic, not a pass/fail gate.
\item Freeze a practical-equivalence margin for B0-A near-null backups before adaptation.
\end{itemize}

\subsection{P3 --- Recoverability}\label{p3-recoverability}

\begin{itemize}
\tightlist
\item Synthetic integer arithmetic, modular arithmetic, Boolean/symbolic logic.
\item Calibrate dev difficulty before probe fitting; report difficulty-stratified results.
\item Natural trajectory collection remains deterministic.
\item I0 primary: exact pre-stop prefix, stochastic continuation, EOS allowed.
\item I1 secondary: forced-compute continuation with EOS constrained.
\item I2 robustness: fixed procedural re-check instruction with no new task facts.
\item Do not pool I0/I1/I2 into one recoverability label.
\item K is REOPENED. Simulate K in {8,16,32} and freeze by a pre-specified precision/cost rule before P3.
\item Primary predictive score: held-out continuation-level Bernoulli / hierarchical-binomial negative log likelihood or deviance.
\item Baselines include prompt-only and question-only hidden-state predictors before J-space.
\end{itemize}

\subsection{Stop rules}\label{stop-rules}

If Gate 0 fails, do not run constrained adaptation as a reconstitution experiment. If the mask audit fails, patch under a new protocol version. If conditional t0 backups already explain recovery, classify self-repair rather than reconstitution. If recoverability does not beat simple baselines, deprioritize J-space causal work.


\section{Freeze Register}
\begin{longtable}{@{}p{0.20\textwidth}p{0.15\textwidth}p{0.25\textwidth}p{0.27\textwidth}@{}}
\caption{v0.6 freeze register. REVIEW-REOPENED items must be re-frozen before the corresponding result-bearing stage.}\label{tab:freeze-register}\\
\toprule
\textbf{Item} & \textbf{Status} & \textbf{Value} & \textbf{Change rule} \\
\midrule
\endfirsthead
\toprule
\textbf{Item} & \textbf{Status} & \textbf{Value} & \textbf{Change rule} \\
\midrule
\endhead
Primary H causal phenotype & PILOT-LOCKED & FaithEval Counterfactual Context & new protocol version \\
H identification training set & PILOT-LOCKED & TriviaQA & new protocol version \\
OOD diagnostics & PILOT-LOCKED & NQ-Open; BioASQ; NonExist diagnostic & new protocol version \\
CETT thresholds & PILOT-LOCKED & 0.0005;0.001;0.002 primary 0.001 & new protocol version \\
FaithEval H effect gate & PILOT-LOCKED & >=5 pp reduction vs baseline & compute-screening gate only \\
Matched-control selectivity & PILOT-LOCKED & >=3 pp; bootstrap CI excludes 0 & compute-screening gate only \\
PPL damage boundary & PILOT-LOCKED & <10\% relative increase & new protocol version \\
Mask tolerance & PILOT-LOCKED & abs <=1e-7 & ledgered bug fix only \\
Sparse-probe full pipeline & PILOT-LOCKED v0.6 & pinned upstream commit + exact source grid/solver/scaling/tolerance & new protocol version \\
Probe stability reporting & PILOT-LOCKED v0.6 & >=10 resampled fits where feasible; selection frequency/Jaccard + causal-mask reproducibility & no arbitrary Jaccard cutoff \\
B0 equivalence rule & PILOT-LOCKED v0.6 & delta=min(2 pp, 20\% of R1 loss); paired 90\% CI; >=80\% power required & new protocol version \\
Conditional backup method & PILOT-LOCKED v0.6 & CoAx + reference validation + target candidate-pool expansion sensitivity & new protocol version \\
Recoverability primary regime & PILOT-LOCKED v0.6 & I0 natural-support, EOS allowed & new protocol version \\
Termination-barrier control & PILOT-LOCKED v0.6 & I0b: force one non-EOS token, then re-enable EOS & new protocol version \\
Forced/procedural regimes & PILOT-LOCKED v0.6 & I1 secondary forced compute; I2 procedural robustness & new protocol version \\
Recoverability K & TO FREEZE & smallest K meeting pre-specified NLL Monte Carlo SE criterion in 8/16/32, else expand & freeze before P3 \\
Natural decoding & PILOT-LOCKED & temperature 0 & new protocol version \\
Recoverability metric & PILOT-LOCKED v0.6 & held-out Bernoulli/hierarchical-binomial NLL/deviance & new protocol version \\
J-space screen & REVIEW-REOPENED & TBD after revised baseline benchmark & freeze before J-space claim \\
Full adaptation optimizer & MAIN-RUN-TBD & TBD & freeze after benchmark + human review \\
Full adaptation LR & MAIN-RUN-TBD & TBD & freeze after benchmark + human review \\
Full adaptation token budget & MAIN-RUN-TBD & TBD & freeze after benchmark + human review \\
Checkpoint token counts & MAIN-RUN-TBD & log-spaced rule; exact counts TBD & freeze after benchmark \\
Full adaptation seeds & MAIN-RUN-TBD & TBD & freeze after pilot variance \\
R2/R3/R4 thresholds & MAIN-RUN-TBD & TBD & freeze before main run \\
\bottomrule
\end{longtable}


\section{Decision Ledger}
\begin{itemize}
\tightlist
\item D-001: Adaptive reconstitution requires learning; fixed-weight repeated probing is insufficient.
\item D-002: Time-zero redundancy mapping is mandatory before adaptation.
\item D-003: R0 establishes selective causal contribution, not mechanism identity.
\item D-004: FaithEval Counterfactual Context remains the pilot primary H-Neuron causal phenotype.
\item D-005: LoRA-only recovery is compensation in the permitted subspace unless later broad adaptation supports more.
\item D-006: v0.5 reopened recoverability K; v0.6 adds the freeze criterion; K=8 is no longer pilot-locked.
\item D-007: Natural-support continuation with EOS allowed becomes the primary recoverability construct.
\item D-008: Hard EOS blocking is retained only as a secondary forced-compute intervention.
\item D-009: Procedural re-check is a separate robustness regime and is not pooled with I0/I1.
\item D-010: Conditional backup discovery is added to t0; marginal intact-model importance is insufficient when self-repair is possible.
\item D-011: Residual-stream decodability is diagnostic, not a hard gate.
\item D-012: Weak pre-existing backup plus training-dependent causal growth is labelled adaptive amplification, not de novo reconstitution.
\item D-013: Strong t0 backup is labelled self-repair / pre-existing redundancy.
\item D-014: The strongest label is de novo / near-null functional reconstitution and requires B0-A plus R0--R4 and cross-seed replication.
\item D-015: Recoverability representation discovery, intervention tuning and causal testing use separate data splits.
\item D-016: H<->J coupling requires matched propagation controls; a raw downstream effect is not sufficient.
\item D-017: No silent threshold repair; material design changes require a new version before result-bearing runs.
\item D-018: B0-A margin rule is frozen before adaptation as delta=min(2 pp, 20\% of observed R1 loss), with paired equivalence CI and >=80\% power required.
\item D-019: A bounded CoAx pool must be validated both on a reference self-repair benchmark and by target-model candidate-pool expansion sensitivity; toy recovery alone does not establish target completeness.
\item D-020: I0b forces exactly one non-EOS token then re-enables EOS to isolate immediate termination-policy effects.
\item D-021: Residual-stream decodability remains diagnostic; no AUC threshold alone can reclassify an outcome as redistribution.
\item D-022: H-Neuron identification pins the official source pipeline; no arbitrary fixed C/Jaccard threshold replaces source replication.
\item D-023: Frozen and checkpoint-refit H-Neuron detectors are reported together to expose detector drift.
\item D-024: K is selected by a pre-specified Monte Carlo precision/cost criterion before P3, not post hoc.
\end{itemize}


\section{Implementation Contract}

\textbf{Status:} Normative implementation companion to Working Paper v0.3\\
\textbf{Date:} 21 September 2026

\subsection{Repository blueprint}\label{repository-blueprint}

\begin{lstlisting}
where-does-the-function-go/
+--- pyproject.toml
+--- uv.lock / requirements lock
+--- README.md
+--- configs/
|   +--- p1_h_neuron.yaml
|   +--- p2_mask_t0.yaml
|   `--- p3_recoverability.yaml
+--- schemas/
|   +--- run_manifest.schema.json
|   +--- mask.schema.json
|   +--- metric.schema.json
|   `--- trajectory.schema.json
+--- src/wdfg/
|   +--- cli.py
|   +--- config.py
|   +--- provenance.py
|   +--- datasets/
|   |   +--- h_neuron.py
|   |   `--- synthetic_recoverability.py
|   +--- models/
|   |   +--- load.py
|   |   `--- llama_contract.py
|   +--- source_adapters/
|   |   +--- h_neurons.py
|   |   `--- jacobian_lens.py
|   +--- interventions/
|   |   +--- mlp_mask.py
|   |   `--- activation_scale.py
|   +--- capture/
|   |   +--- residual.py
|   |   `--- logits.py
|   +--- stages/
|   |   +--- p1_replication.py
|   |   +--- p2_time_zero.py
|   |   `--- p3_recoverability.py
|   +--- scoring/
|   |   +--- faith_eval.py
|   |   +--- synthetic.py
|   |   `--- gate0.py
|   `--- io/
|       +--- jsonl.py
|       `--- artifacts.py
`--- tests/
    +--- test_llama_contract.py
    +--- test_mask.py
    +--- test_gate0.py
    +--- test_t0_immutable.py
    +--- test_pre_stop.py
    +--- test_continuation_seed.py
    `--- test_synthetic_parser.py
\end{lstlisting}

\subsection{Llama mask contract}\label{llama-mask-contract}

Expected target module per transformer block:

\passthrough{\lstinline!model.model.layers[i].mlp!}

Expected projections:

\begin{itemize}
\tightlist
\item
  \passthrough{\lstinline!gate\_proj!}
\item
  \passthrough{\lstinline!up\_proj!}
\item
  \passthrough{\lstinline!down\_proj!}
\end{itemize}

Expected conceptual forward:

\begin{lstlisting}[language=Python]
u = act_fn(gate_proj(x)) * up_proj(x)
y = down_proj(u)
\end{lstlisting}

Intervention point: input to \passthrough{\lstinline!down\_proj!}.

Preferred mechanism:

\begin{lstlisting}[language=Python]
handle = mlp.down_proj.register_forward_pre_hook(mask_pre_hook)
\end{lstlisting}

Conceptual hook:

\begin{lstlisting}[language=Python]
def mask_pre_hook(module, args):
    (u,) = args
    mask = registry.mask_for_layer(layer_idx, device=u.device, dtype=u.dtype)
    return (u * mask,)
\end{lstlisting}

The actual code must handle the library's exact hook signature and verify shape.

Mask tensor shape should broadcast to:

\passthrough{\lstinline![batch, sequence, intermediate\_size]!}

Recommended stored mask shape:

\passthrough{\lstinline![intermediate\_size]!}

Do not mutate \passthrough{\lstinline!down\_proj.weight!} to impose the persistent constraint.

\subsection{Model contract checks}\label{model-contract-checks}

At startup:

\begin{itemize}
\tightlist
\item
  resolve \passthrough{\lstinline!intermediate\_size!};
\item
  verify each requested neuron index is in range;
\item
  verify target layer exists;
\item
  verify the MLP exposes the expected projections;
\item
  run a one-example identity-mask test;
\item
  abort if contract fails.
\end{itemize}

Pilot runs disable \passthrough{\lstinline!torch.compile!}.

\subsection{Source adapter rule}\label{source-adapter-rule}

P1 H-Neuron identification calls a pinned upstream implementation where possible.

Adapter outputs normalized records:

\begin{lstlisting}
example_id
split
layer_index
neuron_index
cett_score
classifier_weight
is_h_neuron
source_commit
\end{lstlisting}

A local optimized implementation is not allowed to replace the source adapter until equivalence fixtures pass.

\subsection{Run commands}\label{run-commands}

Target CLI surface:

\begin{lstlisting}
wdfg p1 identify --config configs/p1_h_neuron.yaml
wdfg p1 causal   --config configs/p1_h_neuron.yaml --parent-run <id>
wdfg p2 t0       --config configs/p2_mask_t0.yaml --mask <mask.json>
wdfg p3 collect  --config configs/p3_recoverability.yaml
wdfg p3 probe    --config configs/p3_recoverability.yaml --parent-run <id>
wdfg validate-run <run_dir>
\end{lstlisting}

Commands are illustrative API requirements; exact CLI framework is implementation choice.

\subsection{No implicit defaults for scientific parameters}\label{no-implicit-defaults-for-scientific-parameters}

The program may have software defaults for paths.

It must not silently default any scientific parameter that affects the protocol.

Required config keys must fail validation if omitted.

Examples:

\begin{itemize}
\tightlist
\item
  model revision;
\item
  seed;
\item
  coefficient threshold;
\item
  intervention alpha;
\item
  sample IDs/fingerprint;
\item
  K;
\item
  B;
\item
  decoding parameters.
\end{itemize}

\subsection{P1 outputs}\label{p1-outputs}

Required:

\begin{lstlisting}
artifacts/h_neurons.json
artifacts/candidate_scores.parquet
artifacts/faith_eval_predictions.parquet
artifacts/control_predictions.parquet
artifacts/gate0_metrics.json
\end{lstlisting}

\passthrough{\lstinline!h\_neurons.json!} is converted to a canonical mask artifact only after Gate 0 evaluation.

\subsection{P2 outputs}\label{p2-outputs}

Required:

\begin{lstlisting}
artifacts/mask.json
artifacts/mask_audit.json
artifacts/t0_candidate_loci.parquet
artifacts/t0_causal_effects.parquet
\end{lstlisting}

\passthrough{\lstinline!t0\_*!} files are write-once.

The software should refuse overwrite unless an explicit \passthrough{\lstinline!--development-override!} flag is present. Such a run is automatically marked non-result-bearing.

\subsection{P3 natural trajectory representation}\label{p3-natural-trajectory-representation}

Each natural trajectory record contains:

\begin{itemize}
\tightlist
\item
  example ID;
\item
  task family;
\item
  task parameters;
\item
  prompt hash;
\item
  prompt token IDs;
\item
  natural output token IDs;
\item
  parsed natural answer;
\item
  correctness;
\item
  pre-stop prefix token IDs;
\item
  EOS token ID;
\item
  stop-token logit/probability;
\item
  top-1 token ID;
\item
  entropy;
\item
  residual artifact pointer;
\item
  manifest run ID.
\end{itemize}

Large residual tensors are stored outside JSONL.

\subsection{P3 continuation representation}\label{p3-continuation-representation}

Each continuation record contains:

\begin{itemize}
\tightlist
\item
  parent example ID;
\item
  intervention ID (\passthrough{\lstinline!I0!}, \passthrough{\lstinline!I1!}, or \passthrough{\lstinline!I2!});
\item
  k index;
\item
  deterministic sampling seed;
\item
  generated token IDs;
\item
  generated text artifact pointer or text;
\item
  last post-continuation \passthrough{\lstinline!FINAL:!} parse;
\item
  rescue boolean;
\item
  parse status.
\end{itemize}

\subsection{P3 continuation regimes}\label{p3-continuation-regimes}

The implementation must preserve three distinct continuation regimes:

\begin{itemize}
\tightlist
\item \passthrough{\lstinline!I0!}: natural-support continuation from the exact pre-stop prefix, EOS allowed.
\item \passthrough{\lstinline!I1!}: forced-compute continuation, EOS constrained for the configured window.
\item \passthrough{\lstinline!I2!}: procedural re-check with the fixed versioned instruction and no new task facts.
\end{itemize}

The run manifest records the regime explicitly. Analysis code must not pool them into one rescue probability by default.

The continuation count \passthrough{\lstinline!K!} is not hard-coded in v0.6; config validation requires an explicitly frozen value from the post-review precision/cost register.

\subsection{Synthetic generator contract}\label{synthetic-generator-contract}

Generator functions must return:

\begin{lstlisting}[language=Python]
{
  "example_id": "...",
  "family": "integer_arithmetic",
  "parameters": {...},
  "prompt": "... FINAL: <answer> ...",
  "ground_truth": "...",
  "checker_version": "..."
}
\end{lstlisting}

Task difficulty parameters are part of the dataset fingerprint.

Train/val/test generator seeds must not overlap.

\subsection{Dataset fingerprint}\label{dataset-fingerprint}

For each dataset/split, fingerprint:

\begin{itemize}
\tightlist
\item
  canonical list of example IDs;
\item
  normalized prompt text;
\item
  task parameters if generated;
\item
  ground-truth labels;
\item
  preprocessing version.
\end{itemize}

Store SHA-256 of the canonical JSON serialization.

\subsection{Gate 0 evaluator}\label{gate-0-evaluator}

Pure function inputs:

\begin{itemize}
\tightlist
\item
  baseline FaithEval predictions;
\item
  H hard suppression predictions;
\item
  contribution-matched control predictions;
\item
  PPL metrics;
\item
  held-out classifier metrics.
\end{itemize}

Output:

\begin{lstlisting}
{
  "G0_A": "PASS|FAIL",
  "G0_B": "PASS|FAIL",
  "G0_C": "PASS|FAIL",
  "G0_D": "PASS|FAIL",
  "overall": "PASS|FAIL|AMBIGUOUS",
  "evidence": {...}
}
\end{lstlisting}

No manual override in a result-bearing run.

\subsection{Integrity checks for pre-stop recomputation}\label{integrity-checks-for-pre-stop-recomputation}

Natural-run pre-stop record stores:

\begin{itemize}
\tightlist
\item
  pre-stop prefix;
\item
  original top-1 token;
\item
  original EOS probability;
\item
  original logits hash if retained.
\end{itemize}

Recompute from prefix.

Pass if:

\begin{itemize}
\tightlist
\item
  top-1 token is identical;
\item
  absolute EOS-probability difference \textless= configured tolerance.
\end{itemize}

Suggested pilot tolerance: \passthrough{\lstinline!1e-4!} after probability computation in float32.

If fail:

\begin{itemize}
\tightlist
\item
  record \passthrough{\lstinline!FAIL\_DATA\_INTEGRITY!} for that example;
\item
  exclude from recoverability labels;
\item
  do not silently continue.
\end{itemize}

\subsection{Hidden-state capture}\label{hidden-state-capture}

Default:

\begin{itemize}
\tightlist
\item
  last-prefix-token residual state at every transformer layer;
\item
  cast/storage dtype recorded;
\item
  stored using \passthrough{\lstinline!safetensors!} or another deterministic tensor format;
\item
  index file maps example ID -\textgreater{} tensor artifact.
\end{itemize}

Do not serialize arbitrary Python pickle objects for scientific artifacts.

\subsection{J-lens adapter}\label{j-lens-adapter}

J-lens is optional in first P3 pass.

Adapter must not change the underlying P3 trajectory collection.

Feature extraction runs as a downstream analysis against frozen trajectory IDs.

This ensures a failed J-lens integration cannot alter the baseline recoverability dataset.

\subsection{Concurrency}\label{concurrency}

Parallelism must not change random results.

Continuation seed derivation uses example/intervention/k identity, not global RNG order.

All stage outputs are sorted by stable example ID before hashing.

\subsection{Error handling}\label{error-handling}

Scientific failures and software failures are separate.

Examples:

\begin{itemize}
\tightlist
\item
  Gate 0 fail = scientific/protocol result.
\item
  Missing mask coordinate = implementation failure.
\item
  Dataset hash mismatch = data-integrity failure.
\item
  CUDA OOM = hardware failure.
\end{itemize}

Only the first category contributes directly to scientific interpretation.

\subsection{Security / privacy}\label{security-privacy}

Pilot datasets should contain no personal or sensitive user data.

Prompt artifacts should be inspectable and releasable where licensing permits.

\subsection{Acceptance criterion for implementation v0.1}\label{acceptance-criterion-for-implementation-v0.1}

The implementation is ready for P1 result-bearing execution when:

\begin{itemize}
\tightlist
\item
  schema validation passes;
\item
  unit tests pass;
\item
  identity mask matches baseline;
\item
  targeted mask zeros intended coordinates;
\item
  H-Neuron source adapter fixture passes;
\item
  dataset fingerprint is stable across reruns;
\item
  run manifest is complete;
\item
  Gate 0 evaluator passes fixture tests;
\item
  dirty worktree is false.
\end{itemize}


\clearpage
\printbibliography[heading=bibintoc,title={References}]
\end{document}
